\chapter{Granuloma Annulare}
\index{Granuloma Annulare}
\label{sec:Granuloma Annulare}

\section{Overview}
Granuloma annulare (GA) is a benign, chronic inflammatory skin condition characterized by annular (ring-shaped) papules or plaques, most commonly on the extremities. It is a self-limited dermatosis of unknown cause that affects people of all ages, with a predilection for children and young adults. The condition is not contagious and typically resolves spontaneously over months to years.
\section{Classification}

\begin{description}[font=\normalfont\bfseries,leftmargin=3em,labelwidth=2.5em]
  \item[Cause] Idiopathic / Inflammatory
  \item[Organ System] Skin
  \item[Pathophysiology] Lymphohistiocytic granulomatous inflammation with necrobiotic collagen degeneration
  \item[Duration] Chronic (months to years)
  \item[Onset] Gradual
  \item[Mode of Transmission] Non-communicable
  \item[Clinical Course] Chronic, self-limited
  \item[Severity] Mild
  \item[Extent of Disease] Localized to Generalized
  \item[Age Group] Children and Young Adults (peak <30 years)
  \item[Epidemiology] Prevalence 0.1--0.4\%; female predominance 2:1
  \item[ICD-11 Code] EE71.0
\end{description}

\section{Causes}
The exact cause of granuloma annulare is unknown. The condition is thought to represent a delayed-type hypersensitivity reaction with T-cell-mediated immune response. Associations have been reported with diabetes mellitus, thyroid disease, dyslipidemia, and malignancy (particularly in generalized GA). Trauma, sun exposure, and viral infections have been proposed as triggers in susceptible individuals.

\section{Risk Factors}

\begin{itemize}[noitemsep]
  \item Female sex
  \item Diabetes mellitus (particularly generalized GA)
  \item Thyroid disease
  \item Dyslipidemia
  \item Malignancy (generalized GA in older adults)
  \item HIV infection
  \item Family history (rare)
  \item UV exposure (possibly)
\end{itemize}

\section{Signs and Symptoms}

\subsection{Early symptoms}
\begin{itemize}[noitemsep]
  \item Early manifestations of granuloma annulare may be subtle and nonspecific
\end{itemize}

\subsection{Common symptoms}
\begin{itemize}[noitemsep]
  \item \subsection{Common symptoms}
  \item \textbf{Common symptoms:}
  \item Firm, smooth, flesh-colored to erythematous papules arranged in an annular (ring) configuration
  \item Most commonly on the dorsal hands, feet, wrists, and ankles
  \item Central clearing with raised border (classic presentation)
  \item Generally asymptomatic; mild pruritus in some cases
  \item No scale, no ulceration
  \item Variants: localized (most common), generalized, subcutaneous, perforating, patch
  \item Difficulty performing daily activities (cosmetic concern)
\end{itemize}

\subsection{Less common symptoms}
\begin{itemize}[noitemsep]
  \item Atypical or less common presentations of granuloma annulare may occur
\end{itemize}

\subsection{Emergency warning signs}
\begin{itemize}[noitemsep]
  \item Severe or worsening symptoms of granuloma annulare require immediate medical evaluation
\end{itemize}
\section{Progression and Stages}
Localized GA typically persists for 1--2 years before spontaneous resolution, though it may last longer. Generalized GA tends to have a more chronic course lasting 3--5 years or more. Recurrence occurs in up to 40\% of cases. Subcutaneous GA (common in children) usually resolves within 1--2 years. The condition does not cause permanent scarring.

\section{Complications}

\begin{itemize}[noitemsep]
  \item Cosmetic concern and psychosocial impact
  \item Large or disseminated lesions
  \item Association with underlying systemic disease (diabetes, malignancy in generalized form)
  \item Recurrence after resolution
  \item Reduced quality of life
  \item Emotional and psychological effects
\end{itemize}

\section{Diagnosis}

\begin{itemize}[noitemsep]
  \item Clinical diagnosis based on characteristic annular morphology and distribution
  \item Skin biopsy shows palisading granulomas with necrobiotic collagen and mucin deposition
  \item Biopsy may be needed to differentiate from tinea corporis, sarcoidosis, erythema migrans
  \item Blood glucose or HbA1c testing to screen for diabetes (generalized GA)
  \item Lipid panel and thyroid function tests as clinically indicated
\end{itemize}

\section{Prevention}

\begin{itemize}[noitemsep]
  \item No established prevention measures
  \item Manage underlying conditions (diabetes, thyroid disease)
  \item Sun protection may help in photosensitive cases
  \item Go for regular check-ups so any problems are caught early
\end{itemize}

\section{Treatment Options}

\begin{itemize}[noitemsep]
  \item Observation and reassurance (many cases resolve spontaneously)
  \item High-potency topical corticosteroids under occlusion
  \item Intralesional corticosteroid injections for localized lesions
  \item Topical calcineurin inhibitors (tacrolimus, pimecrolimus)
  \item Phototherapy (PUVA, narrowband UVB) for generalized GA
  \item Systemic therapies for refractory generalized GA (hydroxychloroquine, dapsone, methotrexate, TNF inhibitors)
  \item Cryotherapy for small lesions
  \item Lifestyle changes and self-care
\end{itemize}

\section{Living with It}

\begin{itemize}[noitemsep]
  \item Follow your treatment plan as prescribed
  \item Keep up with regular appointments with your healthcare provider
  \item Adjust your daily habits as needed
  \item Reach out to family, friends, or support groups for help
  \item Keep learning about your condition and how to manage it
\end{itemize}

\section{Outlook}
Granuloma annulare is a benign, self-limited condition with an excellent prognosis. Approximately 50\% of localized cases resolve within 2 years. Generalized GA tends to persist longer but eventually resolves in most cases. Treatment accelerates clearance but does not prevent recurrence. The condition does not cause permanent tissue damage or increase mortality, though underlying disease associations should be evaluated in generalized cases.
